% Generated by Listas/make_statements.py from pratica-lecture-05.md, item 15. Do not edit.
\textbf{Source.} Muñoz-Garcia \& Toro-Gonzalez (2019), \emph{Strategy and Game Theory: Practice Exercises with Answers}, Ch. 6, Exercise 9 ``Representing Feasible, Individually Rational Payoffs'', p.~305 {[}PDF 316{]}. Answer in the source: yes, a worked answer follows the exercise. Reproduced verbatim (statement only; the book's worked answer follows it).



\begin{quote}

Consider an infinitely-repeated game where the stage game is depicted Table 6.6. Player 1 chooses rows and player 2 chooses columns. Players discount the future using a common discount factor \(\delta\).



\textbf{Part (a)} What outcomes in the stage-game are consistent with Nash equilibrium play?



\textbf{Part (b)} Let \(v_1\) and \(v_2\) be the repeated game payoffs to player 1 and player 2 respectively. Draw the set of feasible payoffs from the repeated game, explaining any normalization you use.



\textbf{Part (c)} Find the set of individually rational feasible set of payoffs.



\textbf{Part (d)} Find a Nash equilibrium in which the players obtain (9, 9) each period. What restrictions on \(\delta\) are necessary?



{\def\LTcaptype{none} % do not increment counter

\begin{longtable}[]{@{}lll@{}}

\toprule\noalign{}

Player 1 ~Player 2 & \(L\) & \(R\) \\

\midrule\noalign{}

\endhead

\bottomrule\noalign{}

\endlastfoot

\(U\) & 9, 9 & 1, 10 \\

\(D\) & 10, 1 & 7, 7 \\

\end{longtable}

}



\emph{Table 6.6} Simultaneous-move game

\end{quote}

